\begin{abcliste}
\abc Photons have no mass, but they have momentum, $p = \frac{h}{\lambda} = \frac{E}{c}$:
\begin{align}
 p &= \pgF = \frac{\nch}{\la} = \pg \approx \result{\pgP{0}{2}}
\end{align}

\abc $N = \frac{P}{E}$ photons per second each give up their momentum $p = \frac{E}{c}$ to the surface. The force is the momentum delivered per second:
\begin{align}
 F &= N\cdot p = \frac{P}{E}\cdot\frac{E}{c} = \FgF \\
   &= \frac{\Pg}{\ncc} \\
   &= \Fg \approx \result{\FgP{-12}{2}}
\end{align}
Tiny, but measurable. (A mirror would feel twice the force, since it reverses the momentum of each photon.)

\abc $p = \frac{h}{\lambda}$: the shorter wavelength gives $\result{\text{more momentum}}$:
\begin{align}
 p &= \pvF = \frac{\nch}{\lb} = \pv \approx \pvP{0}{2}
\end{align}
\end{abcliste}
